\begin{table}[h!]
    \scriptsize
    \captionsetup{width=\linewidth}
    \Caption{\label{tab:ranks_medios_com_pli} Ranks medios por metodo, incluindo o PLIB como referencia complementar. Legenda: AT representa ACO\_R+TS e HR representa HHO+RVNS; Metodo indica a abordagem avaliada; $\bar{r}$ e o rank medio; $\tilde{r}$ e o rank mediano; $M/E$ indica melhor valor ou empate. Em cada instancia, os metodos sao ordenados pelo valor obtido, com menor valor recebendo menor rank; portanto, menor $\bar{r}$ indica melhor desempenho relativo. Esta tabela deve ser lida em conjunto com o teste de Friedman, pois os ranks descrevem a ordem media dos metodos, enquanto o valor-p indica se as diferencas globais sao estatisticamente significativas.}
    \IBGEtab{}{
        \begin{tabular}{lllll}
            \toprule
            Conj. & Metodo & $\bar{r}$ & $\tilde{r}$ & $M/E$ \\
            \midrule \midrule
            TODOS & PLIB & 1.623 & 1.5 & 101 \\
            TODOS & AT & 2.054 & 2 & 64 \\
            TODOS & HR & 2.323 & 2 & 45 \\
            CUBIC & AT & 1.117 & 1 & 28 \\
            CUBIC & PLIB & 2.183 & 2 & 5 \\
            CUBIC & HR & 2.7 & 3 & 0 \\
            DIMACS & PLIB & 1.52 & 1.5 & 49 \\
            DIMACS & HR & 1.99 & 2 & 30 \\
            DIMACS & AT & 2.49 & 2.5 & 19 \\
            HB & PLIB & 1.39 & 1 & 47 \\
            HB & AT & 2.18 & 2 & 17 \\
            HB & HR & 2.43 & 2.5 & 15 \\
            \bottomrule
        \end{tabular}
    }{
    \Fonte{Autor (2026).}
}
\end{table}
